\chapter{Command line and AI assistants}

\file{modbench-cli} reads devices and works with maps without the window: for
scripts, for servers without a screen, and for a quick check. It comes with
the app. It never writes to a device.

\section{Commands}

\begin{lstlisting}[language=bash]
modbench-cli read --tcp 192.168.1.10 holding 0x0405 1
modbench-cli read --udp 192.168.1.10 holding 0 10
modbench-cli map --rtu COM3 --unit 18 device.mbmap.json --filter voltage
modbench-cli scan --rtu COM3 --baud 9600 --from 1 --to 32
modbench-cli id --tcp 192.168.1.10
modbench-cli simulate device.mbmap.json --strict
modbench-cli simulate device.mbmap.json --generators moving.json
modbench-cli simulate device.mbmap.json --rtu-over-tcp
modbench-cli simulate device.mbmap.json --rtu COM5 --baud 19200 --parity E
modbench-cli validate device.mbmap.json
modbench-cli convert old-device.mbmap
\end{lstlisting}

\begin{center}
\begin{tabularx}{\linewidth}{l L}
  \toprule
  \thead{Command} & \thead{What it does} \\
  \midrule
  \val{read} & Reads coils or registers, like a manual request. \\
  \val{map} & Reads every register of a map from a device, by name;
    \val{--filter} keeps the names that match. \\
  \val{scan} & Finds the unit IDs that answer. \\
  \val{id} & Device identification (FC43). \\
  \val{simulate} & A simulated device shaped after a map (port 5020 by
    default); \val{--strict} answers only the map's addresses. \\
  \val{validate} & Checks a map against the format, with the same messages as
    the app. \\
  \val{convert} & Converts a version 1 map to a \file{.mbmap.json} next to it. \\
  \bottomrule
\end{tabularx}
\end{center}

\val{modbench-cli --help} lists every command and option.

\section{Generators file}

\val{--generators} takes a JSON list. Each entry names a map register (in
engineering units) or a raw \val{area} and \val{address}, a \val{kind}, and how
often it changes (\val{every\_ms}, 1000 by default, at least 50):

\begin{lstlisting}
[
  {"register": "Voltage L1", "kind": "sine", "offset": 230, "amplitude": 5, "period_s": 10, "every_ms": 500},
  {"area": "input", "address": 0, "kind": "increment", "step": 1, "min": 0, "max": 100},
  {"area": "holding", "address": 1, "kind": "random", "min": 10, "max": 20},
  {"area": "discrete", "address": 0, "kind": "copy", "from": {"area": "coil", "address": 0}}
]
\end{lstlisting}

\section{AI assistants (MCP)}

\val{modbench-cli mcp} serves the same reads over the Model Context Protocol
(stdio), so an AI assistant such as Claude Code can read devices and check
maps: \val{read\_map} reads a device with a map, \val{validate\_map} checks a
map file. Like the command line, it never writes to a device.
